\documentclass{article}



\usepackage{arxiv}

\usepackage[utf8]{inputenc} % allow utf-8 input
\usepackage[T1]{fontenc}    % use 8-bit T1 fonts
\usepackage{hyperref}       % hyperlinks
\usepackage{url}            % simple URL typesetting
\usepackage{booktabs}       % professional-quality tables
\usepackage{amsfonts}       % blackboard math symbols
\usepackage{nicefrac}       % compact symbols for 1/2, etc.
\usepackage{microtype}      % microtypography
\usepackage{lipsum}		% Can be removed after putting your text content
\usepackage{graphicx}
\usepackage{natbib}
\usepackage{doi}

\usepackage{tikz}
\usepackage{amsmath}

\DeclareMathOperator{\E}{\mathbb{E}}



\title{Expected number of $n$-sided die throws to collect $k$ points is a geometric series for $k\leq n$}

%\date{September 9, 1985}	% Here you can change the date presented in the paper title
%\date{} 					% Or removing it

\author{ \href{https://orcid.org/0000-0003-2284-3204}{\includegraphics[scale=0.06]{orcid.pdf}\hspace{1mm}Boris~Demeshev}\\ %\thanks{github.com/bdemeshev/me} \\
	Department of Computer Science, Department of Applied Economics\\
	State University Higher School of Economics\\
	Moscow \\
	\texttt{bdemeshev@hse.ru} \\
	%% examples of more authors
	%% \AND
	%% Coauthor \\
	%% Affiliation \\
	%% Address \\
	%% \texttt{email} \\
	%% \And
	%% Coauthor \\
	%% Affiliation \\
	%% Address \\
	%% \texttt{email} \\
	%% \And
	%% Coauthor \\
	%% Affiliation \\
	%% Address \\
	%% \texttt{email} \\
}

% Uncomment to remove the date
%\date{}

% Uncomment to override  the `A preprint' in the header
%\renewcommand{\headeright}{Technical Report}
%\renewcommand{\undertitle}{Technical Report}
\renewcommand{\shorttitle}{Expected number of throws}

%%% Add PDF metadata to help others organize their library
%%% Once the PDF is generated, you can check the metadata with
%%% $ pdfinfo template.pdf
\hypersetup{
pdftitle={A template for the arxiv style},
pdfsubject={q-bio.NC, q-bio.QM},
pdfauthor={Boris~Demeshev},
pdfkeywords={First keyword, Second keyword, More},
}

\begin{document}
\maketitle

% \begin{abstract}
%	Expected number of $n$-sided die throws to collect $k$ points 
%	is a geometric series for $k \in \{1, 2 \ldots, n\}$.
% \end{abstract}


% keywords can be removed
% \keywords{First keyword \and Second keyword \and More}


\section{Link between die and geometric series}
Let's start with a visual proof of a formula for the finite geometric series sum.
We consider a growing geometric series $(b_j)$ where $b_{j + 1} = r b_j$ with ratio $r > 1$.
We represent the ratio as $r = 1 + \alpha$, so $b_{j + 1} = b_j + \alpha b_j$.


\tikzset{ext/.pic={
\path [fill=white] (-0.2,0)to[bend left](0,0.1)to[bend right](0.2,0.2)to(0.2,0)to[bend left](0,-0.1)to[bend right](-0.2,-0.2)--cycle;
\draw (-0.2,0)to[bend left](0,0.1)to[bend right](0.2,0.2) (0.2,0)to[bend left](0,-0.1)to[bend right](-0.2,-0.2);
}}

%\begin{tikzpicture}
%\foreach \Point/\PointLabel in {(0,0)/0, (1,0)/b_1, (1.8,0)/b_2, (2.9,0)/b_3, (4.9,0)/b_4, (8.4,0)/b_{k-1}, (14.2,0)/b_k}
%\draw[fill=black] \Point circle (0.05) node[below] {$\PointLabel$};
%\path [draw] (0,0)--pic [rotate=90]{ext}(13,0);
%\path [draw] (13,0) -- (16,0);
%\end{tikzpicture}

\begin{tikzpicture}
\foreach \Point/\PointLabel in {(0,0)/0, (1,0)/b_1, (1.8,0)/b_2, (2.9,0)/b_3, (4.9,0)/b_4, (8.4,0)/b_{k-1}, (14.2,0)/b_k}
\draw[fill=black] \Point circle (0.05) node[below] {$\PointLabel$};
\path [draw] (0,0)--pic [rotate=90]{ext}(13,0);
\path [draw] (13,0) -- (16,0);

\foreach \Point/\PointLabel in {(0.45,0)/b_1, (1.4,0)/\alpha b_1, (2.35,0)/\alpha b_2, (3.9,0)/\alpha b_3, (11.5,0)/\alpha b_{k-1}}
\draw \Point node[above] {$\PointLabel$};
\end{tikzpicture}
		

Adding all segments lengths we see that
\[
b_{k} = b_1 + \alpha (b_1 + b_2 + b_3 + \ldots + b_{k-1})
\]

This property of geometric series also has an economic intuition. 
You final wellfare $b_k$ equals your initial wellfare $b_1$ plus all interest payments,
where $\alpha$ is the nominal rate in the compound interest scheme.

% check! proofs without words! ?

Now let's toss a $n$-sided die until we collect $k$ points or more with $k\leq n$.
We denote the random number of tosses by $X_k$ and the expected number of tosses 
by $b_k = \E(X_k)$. 

For sure we'll collect one point or more with exactly one toss, $X_1 = 1$ and hence $b_1 = 1$. 

Now we need to collect $k$ points or more. Let's consider the first toss.
With probability $(n - k)/n$ zero more tosses will be required. 
And with probability $1/n$ we will need to collect $1$, $2$, \ldots, $k-1$ points.

The recurrence formula is
\[
b_k = 1 + \frac{n - k}{n} \cdot 0 + \frac{1}{n} (b_1 + b_2 + \ldots + b_{k-1}).
\]
This equation coupled with initial condition defines the geometric series with $b_1 = 1$ and ratio $r = 1 + \frac{1}{n}$.

Hence we obtain
\[
b_k = \E(X_k) = \left( 1 + \frac{1}{n} \right)^{k-1}
\]


This result appears in disguised form $\E(X_k) = \sum_{i=0}^{k-1} \binom{k-1}{i} / n^i$ in \cite{conroy2021collection}.
The particular case with $k=n$ is proven in \cite{trevino2020expected}.
In both these sources the Christmas stockings theorem is invoked in the proof. 
The approach with recurrence equation is used in \cite{146114}, 
where the chase for the general $k$ hides the simple structure for $k\leq n$. 

\bibliographystyle{unsrtnat}
\bibliography{references}  %%% Uncomment this line and comment out the ``thebibliography'' section below to use the external .bib file (using bibtex) .


%%% Uncomment this section and comment out the \bibliography{references} line above to use inline references.
% \begin{thebibliography}{1}

% 	\bibitem{kour2014real}
% 	George Kour and Raid Saabne.
% 	\newblock Real-time segmentation of on-line handwritten arabic script.
% 	\newblock In {\em Frontiers in Handwriting Recognition (ICFHR), 2014 14th
% 			International Conference on}, pages 417--422. IEEE, 2014.

% 	\bibitem{kour2014fast}
% 	George Kour and Raid Saabne.
% 	\newblock Fast classification of handwritten on-line arabic characters.
% 	\newblock In {\em Soft Computing and Pattern Recognition (SoCPaR), 2014 6th
% 			International Conference of}, pages 312--318. IEEE, 2014.

% 	\bibitem{hadash2018estimate}
% 	Guy Hadash, Einat Kermany, Boaz Carmeli, Ofer Lavi, George Kour, and Alon
% 	Jacovi.
% 	\newblock Estimate and replace: A novel approach to integrating deep neural
% 	networks with existing applications.
% 	\newblock {\em arXiv preprint arXiv:1804.09028}, 2018.

% \end{thebibliography}


\end{document}
